\subsection{Upper and Lower Bounds}

Bounds are the first level of extremal information.  They compare one
candidate element with every element of a nonempty set, but they do not
require the candidate to be sharp or to belong to the set.

\begin{figure}[htbp]
  \centering
  \begin{tikzpicture}[
  x=1cm,
  y=1cm,
  >=Latex,
  axis/.style={thick},
  tick/.style={thick},
  pt/.style={circle, fill=black, inner sep=1.7pt},
  cand/.style={circle, fill=blue!60!black, inner sep=2pt},
  witness/.style={circle, fill=red!70!black, inner sep=2pt},
  lbl/.style={font=\small},
  note/.style={font=\footnotesize,text=gray!70!black}
]

\foreach \y/\title in {3/{upper bound},1/{not upper bound},-1/{lower bound},-3/{not lower bound}}{
  \draw[axis] (-3.2,\y) -- (3.2,\y);
  \node[lbl,anchor=east] at (-3.35,\y) {\title};
}

% Upper bound.
\foreach \x in {-2.1,-1.1,0.15,0.85}
  \node[pt] at (\x,3) {};
\node[cand] at (1.9,3) {};
\draw[blue!60!black,thick] (1.9,2.78) -- (1.9,3.22);
\node[lbl,blue!60!black,above] at (1.9,3.22) {$u$};
\node[note,below] at (-0.6,2.78) {$x\leq u$ for all $x\in A$};

% Not an upper bound.
\foreach \x in {-2.1,-1.1,0.15,1.65}
  \node[pt] at (\x,1) {};
\node[cand] at (0.85,1) {};
\node[witness] at (1.65,1) {};
\draw[blue!60!black,thick] (0.85,0.78) -- (0.85,1.22);
\node[lbl,blue!60!black,above] at (0.85,1.22) {$u$};
\node[lbl,red!70!black,above] at (1.65,1.22) {$x$};
\node[note,below] at (1.25,0.78) {$u<x$};

% Lower bound.
\foreach \x in {-0.85,0.15,1.1,2.1}
  \node[pt] at (\x,-1) {};
\node[cand] at (-1.9,-1) {};
\draw[blue!60!black,thick] (-1.9,-1.22) -- (-1.9,-0.78);
\node[lbl,blue!60!black,above] at (-1.9,-0.78) {$\ell$};
\node[note,below] at (0.4,-1.22) {$\ell\leq x$ for all $x\in A$};

% Not a lower bound.
\foreach \x in {-1.65,-0.15,1.1,2.1}
  \node[pt] at (\x,-3) {};
\node[cand] at (-0.85,-3) {};
\node[witness] at (-1.65,-3) {};
\draw[blue!60!black,thick] (-0.85,-3.22) -- (-0.85,-2.78);
\node[lbl,blue!60!black,above] at (-0.85,-2.78) {$\ell$};
\node[lbl,red!70!black,above] at (-1.65,-2.78) {$x$};
\node[note,below] at (-1.25,-3.22) {$x<\ell$};

\end{tikzpicture}

  \caption{Upper and lower bound conditions on a nonempty set.}
  \label{fig:upper-lower-bound-combinations}
\end{figure}

\begin{definitionbox}{Bound}
\begin{definition}[Bound]
\label{def:bound}
Let $(S,\leq)$ be an ordered set, let $A\subseteq S$ be nonempty, and let
$b\in S$.  A \emph{bound} of $A$ means a one-sided comparison between $b$
and every element of $A$: either $x\leq b$ for all $x\in A$, or $b\leq x$
for all $x\in A$.
\end{definition}
\end{definitionbox}

\begin{remark*}[Standard quantified statement]
\[
  A\neq\varnothing
  \quad\text{and}\quad
  \bigl[(\forall x\in A)(x\leq b)
  \ \text{or}\ 
  (\forall x\in A)(b\leq x)\bigr].
\]
\end{remark*}

\begin{remark*}[Definition predicate reading]
\[
  \operatorname{Bound}_{S,\leq}(b,A)
  :=
  \operatorname{UB}_{S,\leq}(b,A)
  \lor
  \operatorname{LB}_{S,\leq}(b,A).
\]
\end{remark*}

\begin{remark*}[Negated quantified statement]
\[
  A=\varnothing
  \quad\text{or}\quad
  \bigl[(\exists x\in A)(b<x)
  \ \text{and}\ 
  (\exists y\in A)(y<b)\bigr].
\]
\end{remark*}

\begin{remark*}[Negation predicate reading]
\[
  \neg\operatorname{Bound}_{S,\leq}(b,A)
  \Longleftrightarrow
  A=\varnothing
  \lor
  \bigl(\neg\operatorname{UB}_{S,\leq}(b,A)
  \land
  \neg\operatorname{LB}_{S,\leq}(b,A)\bigr).
\]
\end{remark*}

\begin{remark*}[Failure modes]
The statement fails if the set is empty, or if $b$ lies strictly between
two elements of $A$.
\end{remark*}

\begin{remark*}[Failure mode decomposition]
Failure of the upper-bound side means some element of $A$ is strictly
larger than $b$.  Failure of the lower-bound side means some element of
$A$ is strictly smaller than $b$.  Failure to be any bound requires both
witnesses.
\end{remark*}

\begin{dependencies}
\begin{itemize}
  \item \hyperref[def:ordered-set]{Ordered Set}
\end{itemize}
\end{dependencies}

\begin{definitionbox}{Upper Bound}
\begin{definition}[Upper Bound]
\label{def:upper-bound}
Let $(S,\leq)$ be an ordered set, let $A\subseteq S$ be nonempty, and let
$u\in S$.  The element $u$ is an \emph{upper bound} of $A$ if every element
of $A$ is less than or equal to $u$.
\end{definition}
\end{definitionbox}

\begin{remark*}[Standard quantified statement]
\[
  A\neq\varnothing
  \quad\text{and}\quad
  (\forall x\in A)(x\leq u).
\]
\end{remark*}

\begin{remark*}[Definition predicate reading]
\[
  \operatorname{UB}_{S,\leq}(u,A)
  :=
  A\neq\varnothing\land(\forall x\in A)(x\leq u).
\]
\end{remark*}

\begin{remark*}[Negated quantified statement]
\[
  A=\varnothing
  \quad\text{or}\quad
  (\exists x\in A)(u<x).
\]
\end{remark*}

\begin{remark*}[Negation predicate reading]
\[
  \neg\operatorname{UB}_{S,\leq}(u,A)
  \Longleftrightarrow
  A=\varnothing\lor(\exists x\in A)(u<x).
\]
\end{remark*}

\begin{remark*}[Failure modes]
The set may be empty, or a counterexample element of $A$ may sit strictly
above the proposed bound.
\end{remark*}

\begin{remark*}[Failure mode decomposition]
To disprove that $u$ is an upper bound, produce either the structural
failure $A=\varnothing$ or a witness $x\in A$ with $u<x$.
\end{remark*}

\begin{dependencies}
\begin{itemize}
  \item \hyperref[def:bound]{Bound}
\end{itemize}
\end{dependencies}

\begin{definitionbox}{Lower Bound}
\begin{definition}[Lower Bound]
\label{def:lower-bound}
Let $(S,\leq)$ be an ordered set, let $A\subseteq S$ be nonempty, and let
$\ell\in S$.  The element $\ell$ is a \emph{lower bound} of $A$ if every
element of $A$ is greater than or equal to $\ell$.
\end{definition}
\end{definitionbox}

\begin{remark*}[Standard quantified statement]
\[
  A\neq\varnothing
  \quad\text{and}\quad
  (\forall x\in A)(\ell\leq x).
\]
\end{remark*}

\begin{remark*}[Definition predicate reading]
\[
  \operatorname{LB}_{S,\leq}(\ell,A)
  :=
  A\neq\varnothing\land(\forall x\in A)(\ell\leq x).
\]
\end{remark*}

\begin{remark*}[Negated quantified statement]
\[
  A=\varnothing
  \quad\text{or}\quad
  (\exists x\in A)(x<\ell).
\]
\end{remark*}

\begin{remark*}[Negation predicate reading]
\[
  \neg\operatorname{LB}_{S,\leq}(\ell,A)
  \Longleftrightarrow
  A=\varnothing\lor(\exists x\in A)(x<\ell).
\]
\end{remark*}

\begin{remark*}[Failure modes]
The set may be empty, or a counterexample element of $A$ may sit strictly
below the proposed bound.
\end{remark*}

\begin{remark*}[Failure mode decomposition]
To disprove that $\ell$ is a lower bound, produce either the structural
failure $A=\varnothing$ or a witness $x\in A$ with $x<\ell$.
\end{remark*}

\begin{dependencies}
\begin{itemize}
  \item \hyperref[def:bound]{Bound}
\end{itemize}
\end{dependencies}

\begin{definitionbox}{Bounded Above}
\begin{definition}[Bounded Above]
\label{def:bounded-above}
Let $(S,\leq)$ be an ordered set and let $A\subseteq S$ be nonempty.  The
set $A$ is \emph{bounded above} if it has at least one upper bound in $S$.
\end{definition}
\end{definitionbox}

\begin{remark*}[Standard quantified statement]
\[
  A\neq\varnothing
  \quad\text{and}\quad
  (\exists u\in S)(\forall x\in A)(x\leq u).
\]
\end{remark*}

\begin{remark*}[Definition predicate reading]
\[
  \operatorname{BddAbove}_{S,\leq}(A)
  :=
  (\exists u\in S)\operatorname{UB}_{S,\leq}(u,A).
\]
\end{remark*}

\begin{remark*}[Negated quantified statement]
\[
  A=\varnothing
  \quad\text{or}\quad
  (\forall u\in S)(\exists x\in A)(u<x).
\]
\end{remark*}

\begin{remark*}[Negation predicate reading]
\[
  \neg\operatorname{BddAbove}_{S,\leq}(A)
  \Longleftrightarrow
  A=\varnothing
  \lor
  (\forall u\in S)\neg\operatorname{UB}_{S,\leq}(u,A).
\]
\end{remark*}

\begin{remark*}[Failure modes]
The set may be empty, or every proposed upper bound may be defeated by
some larger element of $A$.
\end{remark*}

\begin{remark*}[Failure mode decomposition]
Unboundedness above is not one failed candidate; it is a uniform
procedure that defeats every candidate $u\in S$.
\end{remark*}

\begin{dependencies}
\begin{itemize}
  \item \hyperref[def:upper-bound]{Upper Bound}
\end{itemize}
\end{dependencies}

\begin{definitionbox}{Bounded Below}
\begin{definition}[Bounded Below]
\label{def:bounded-below}
Let $(S,\leq)$ be an ordered set and let $A\subseteq S$ be nonempty.  The
set $A$ is \emph{bounded below} if it has at least one lower bound in $S$.
\end{definition}
\end{definitionbox}

\begin{remark*}[Standard quantified statement]
\[
  A\neq\varnothing
  \quad\text{and}\quad
  (\exists \ell\in S)(\forall x\in A)(\ell\leq x).
\]
\end{remark*}

\begin{remark*}[Definition predicate reading]
\[
  \operatorname{BddBelow}_{S,\leq}(A)
  :=
  (\exists \ell\in S)\operatorname{LB}_{S,\leq}(\ell,A).
\]
\end{remark*}

\begin{remark*}[Negated quantified statement]
\[
  A=\varnothing
  \quad\text{or}\quad
  (\forall \ell\in S)(\exists x\in A)(x<\ell).
\]
\end{remark*}

\begin{remark*}[Negation predicate reading]
\[
  \neg\operatorname{BddBelow}_{S,\leq}(A)
  \Longleftrightarrow
  A=\varnothing
  \lor
  (\forall \ell\in S)\neg\operatorname{LB}_{S,\leq}(\ell,A).
\]
\end{remark*}

\begin{remark*}[Failure modes]
The set may be empty, or every proposed lower bound may be defeated by
some smaller element of $A$.
\end{remark*}

\begin{remark*}[Failure mode decomposition]
Unboundedness below is a uniform defeat of all lower-bound candidates.
\end{remark*}

\begin{dependencies}
\begin{itemize}
  \item \hyperref[def:lower-bound]{Lower Bound}
\end{itemize}
\end{dependencies}

\begin{definitionbox}{Bounded Set}
\begin{definition}[Bounded Set]
\label{def:bounded}
Let $(S,\leq)$ be an ordered set and let $A\subseteq S$ be nonempty.  The
set $A$ is \emph{bounded} if it is bounded above and bounded below.
\end{definition}
\end{definitionbox}

\begin{remark*}[Standard quantified statement]
\[
  A\neq\varnothing
  \quad\text{and}\quad
  (\exists \ell,u\in S)(\forall x\in A)(\ell\leq x\leq u).
\]
\end{remark*}

\begin{remark*}[Definition predicate reading]
\[
  \operatorname{Bdd}_{S,\leq}(A)
  :=
  \operatorname{BddBelow}_{S,\leq}(A)
  \land
  \operatorname{BddAbove}_{S,\leq}(A).
\]
\end{remark*}

\begin{remark*}[Negated quantified statement]
\[
  A=\varnothing
  \quad\text{or}\quad
  (\forall \ell,u\in S)(\exists x\in A)(x<\ell\ \text{or}\ u<x).
\]
\end{remark*}

\begin{remark*}[Negation predicate reading]
\[
  \neg\operatorname{Bdd}_{S,\leq}(A)
  \Longleftrightarrow
  A=\varnothing
  \lor
  \neg\operatorname{BddBelow}_{S,\leq}(A)
  \lor
  \neg\operatorname{BddAbove}_{S,\leq}(A).
\]
\end{remark*}

\begin{remark*}[Failure modes]
The set may be empty, unbounded below, unbounded above, or both.
\end{remark*}

\begin{remark*}[Failure mode decomposition]
Boundedness has two independent one-sided requirements.  A failure on
either side is enough to destroy boundedness.
\end{remark*}

\begin{dependencies}
\begin{itemize}
  \item \hyperref[def:bounded-above]{Bounded Above}
  \item \hyperref[def:bounded-below]{Bounded Below}
\end{itemize}
\end{dependencies}
